\chapter{Per-Request Identity Attribution Specification}\label{ch:identity}

This chapter specifies the per-request identity attribution system, addressing
the MAS MGF requirement REG-MGF-2.1.2-002 that ``each agent action must be
traceable to the requesting identity and approval chain.''

\section{Overview}

Per-request identity attribution ensures that every AI agent action can be
traced to:
\begin{enumerate}
  \item The \textbf{initiating user} who triggered the request
  \item The \textbf{agent identity} that executed the action
  \item The \textbf{approval chain} (if human-in-the-loop was invoked)
  \item The \textbf{system context} (pipeline, model version, timestamp)
\end{enumerate}

\section{Identity model}

\subsection{Agent identity schema}

\begin{lstlisting}[basicstyle=\ttfamily\small,caption={Agent identity record schema.}]
{
  "agent_id": "agent-tb-variance-001",
  "agent_type": "variance_analyzer",
  "model_version": "qwen2.5-72b-instruct-v1.2",
  "pipeline_version": "tb-review-2.1.0",
  "deployment_environment": "production",
  "capabilities": [
    "variance_calculation",
    "anomaly_detection",
    "commentary_generation"
  ],
  "permissions": {
    "can_modify_data": false,
    "can_trigger_downstream": true,
    "requires_human_approval": ["commentary_generation"]
  }
}
\end{lstlisting}

\subsection{Request identity schema}

\begin{lstlisting}[basicstyle=\ttfamily\small,caption={Request identity record schema.}]
{
  "request_id": "req-20260418-001234",
  "correlation_id": "session-abc123",
  "initiating_user": {
    "user_id": "jsmith@company.com",
    "user_type": "human",
    "authentication_method": "oauth2_sso",
    "roles": ["tb_reviewer", "senior_accountant"]
  },
  "timestamp": "2026-04-18T10:30:00+08:00",
  "source_system": "tb-review-dashboard",
  "source_ip": "10.0.1.100"
}
\end{lstlisting}

\section{Audit trail specification}

\subsection{Audit event schema}

Every agent action generates an audit event:

\begin{lstlisting}[basicstyle=\ttfamily\small,caption={Audit event schema.}]
{
  "event_id": "evt-20260418-567890",
  "event_type": "agent_action",
  "timestamp": "2026-04-18T10:30:05+08:00",
  
  "request_identity": {
    "request_id": "req-20260418-001234",
    "initiating_user_id": "jsmith@company.com"
  },
  
  "agent_identity": {
    "agent_id": "agent-tb-variance-001",
    "model_version": "qwen2.5-72b-instruct-v1.2"
  },
  
  "action": {
    "action_type": "variance_calculation",
    "entity_id": "HKG-2026-04",
    "input_hash": "sha256:abc123...",
    "output_hash": "sha256:def456..."
  },
  
  "approval_chain": [
    {
      "approver_id": "jsmith@company.com",
      "approval_type": "implicit",
      "timestamp": "2026-04-18T10:30:00+08:00"
    }
  ],
  
  "outcome": {
    "status": "success",
    "duration_ms": 1234,
    "tokens_used": 500
  }
}
\end{lstlisting}

\subsection{Required audit events}

\begin{table}[htbp]
\centering
\caption{Mandatory audit events by action type.}
\label{tab:audit-events}
\footnotesize
\begin{tabularx}{\linewidth}{@{}l l X@{}}
\toprule
\textbf{Action Type} & \textbf{Event Type} & \textbf{Required Fields} \\ \midrule
AI inference & \texttt{agent\_inference} & model\_version, input\_hash, output\_hash \\
Data modification & \texttt{agent\_write} & entity\_id, before\_hash, after\_hash \\
Human review request & \texttt{review\_requested} & reviewer\_id, queue\_id \\
Human approval & \texttt{review\_approved} & approver\_id, decision, rationale \\
Human rejection & \texttt{review\_rejected} & rejecter\_id, decision, rationale \\
Error/failure & \texttt{agent\_error} & error\_code, stack\_trace\_hash \\
\bottomrule
\end{tabularx}
\end{table}

\section{Storage and retention}

\subsection{Storage requirements}

\begin{table}[htbp]
\centering
\caption{Audit storage requirements.}
\label{tab:audit-storage}
\footnotesize
\begin{tabularx}{\linewidth}{@{}l r X@{}}
\toprule
\textbf{Data Type} & \textbf{Retention} & \textbf{Storage} \\ \midrule
Audit events & 7 years & Append-only log (immutable) \\
Request identities & 7 years & Indexed for query \\
Agent configurations & Forever & Version-controlled \\
Input/output hashes & 7 years & With audit events \\
Full input/output & 90 days & Cold storage, then purge \\
\bottomrule
\end{tabularx}
\end{table}

\subsection{Immutability requirements}

Audit logs \textbf{must} be immutable:
\begin{itemize}
  \item Append-only storage (no updates, no deletes)
  \item Cryptographic chaining (hash of previous event included)
  \item External timestamp authority (prevents backdating)
  \item Write-once media for long-term archive
\end{itemize}

\section{API specification}

\subsection{Audit write API}

Two audit write points are required per action: a
\texttt{pre\_action} reservation (before the inference executes) and
a \texttt{post\_action} completion (after the inference returns). The
\texttt{pre\_action} write is the enforcement gate --- if it fails,
the inference is \emph{rejected before execution}, which is the only
tractable way to achieve the ``no action without audit'' invariant
(see Section~\ref{sec:audit-failure-handling}).

\begin{lstlisting}[language=Python,basicstyle=\ttfamily\small,caption={Audit service API contract.}]
def reserve_agent_action(
    request_identity: RequestIdentity,
    agent_identity: AgentIdentity,
    planned_action: AgentActionPlan,
) -> AuditEventId:
    """
    Reserve an audit event slot BEFORE the agent action executes.

    This function MUST be called and succeed before the inference is
    dispatched. A successful return contractually commits the caller
    to writing a matching complete_agent_action() within the event's
    TTL (default: 120 s).

    Raises:
        AuditWriteError: If the reservation cannot be written. The
            caller MUST NOT dispatch the agent action.
    """

def complete_agent_action(
    event_id: AuditEventId,
    action: AgentAction,
    outcome: ActionOutcome,
) -> None:
    """
    Complete a previously reserved audit event with the actual
    action and outcome. MUST be called after the action returns
    (whether success or error).

    If this call fails, the reservation remains visible in the audit
    trail in state "orphaned" and operations are alerted; the caller
    MUST return an error to the user because downstream side effects
    could not be verifiably attributed.
    """
\end{lstlisting}

\subsection{Audit query API}

\begin{lstlisting}[language=Python,basicstyle=\ttfamily\small,caption={Audit query API contract.}]
def query_audit_trail(
    filters: AuditFilters,
    pagination: Pagination
) -> AuditQueryResult:
    """
    Query the audit trail.
    
    Filters:
        - request_id: Specific request
        - user_id: All actions by user
        - agent_id: All actions by agent
        - entity_id: All actions on entity
        - time_range: Start and end timestamps
        - action_type: Filter by action type
    
    Returns:
        List of matching audit events with pagination
    """
\end{lstlisting}

\section{Integration requirements}

\subsection{Pipeline integration}

Every AI pipeline \textbf{must} integrate the identity attribution system:

\begin{enumerate}
  \item \textbf{Request entry}: Capture initiating user identity from authentication
  \item \textbf{Pre-inference}: Log \texttt{agent\_inference\_start} event
  \item \textbf{Post-inference}: Log \texttt{agent\_inference\_complete} event
  \item \textbf{Data write}: Log \texttt{agent\_write} event before commit
  \item \textbf{Error}: Log \texttt{agent\_error} event on any failure
\end{enumerate}

\subsection{Failure handling}\label{sec:audit-failure-handling}

LLM inference is not reversible once executed: tokens have been
emitted, compute was consumed, and any downstream side effects
(database writes, external calls) are already in flight. ``Roll
back the inference'' is therefore not a valid operation. The
specification enforces audit coverage by \emph{gating} the action
behind a successful reservation write, and by containing side
effects:

\begin{enumerate}
  \item \textbf{Reservation-before-action.} The
        \texttt{reserve\_agent\_action()} call (see audit write
        API) MUST succeed before the agent dispatches the
        inference. If the reservation write fails, the agent
        action MUST NOT be executed and the request MUST return a
        \texttt{503} error with an operations alert.
  \item \textbf{Side-effect containment.} Any data-modification or
        external call triggered by an AI output MUST be staged and
        only committed after \texttt{complete\_agent\_action()}
        succeeds. If completion fails, staged side effects are
        discarded (this \emph{is} a rollback --- of the side
        effects, not of the inference itself).
  \item \textbf{Orphaned reservation handling.} If a reservation
        is written but never completed within the TTL (e.g.\ the
        inference server crashed), the reservation is marked
        \texttt{orphaned}; operations are alerted; staged side
        effects are discarded; downstream consumers see an
        explicit failure.
  \item \textbf{Alerting and no-silent-proceed.} An alert MUST be
        sent to operations on every audit failure; the system
        MUST NOT proceed to emit results without a completed audit
        record.
\end{enumerate}

\noindent
This wording supersedes the earlier text (which called for
``rolling back the agent action''): because LLM inference is not
reversible, the equivalent invariant is enforced by gating with a
pre-action reservation and staging side effects.

\section{Compliance mapping}

\begin{table}[htbp]
\centering
\caption{MAS MGF requirement mapping.}
\label{tab:mgf-mapping}
\footnotesize
\begin{tabularx}{\linewidth}{@{}l l X@{}}
\toprule
\textbf{Requirement} & \textbf{Section} & \textbf{Implementation} \\ \midrule
REG-MGF-2.1.2-002 & \S10.2--10.3 & Request identity schema, audit trail \\
REG-MGF-2.1.3-001 & \S10.4 & Immutable storage, retention policy \\
REG-MGF-2.2.1-001 & \S10.5 & Audit write API with rollback \\
\bottomrule
\end{tabularx}
\end{table}

\section{Implementation checklist and integration status}

\subsection{Audit service (this chapter)}

\begin{itemize}
  \item[\texttt{[x]}] Agent identity schema defined
  \item[\texttt{[x]}] Request identity schema defined
  \item[\texttt{[x]}] Audit event schema defined
  \item[\texttt{[x]}] Audit write API contract specified
        (reserve/complete pattern, Section~\ref{sec:audit-failure-handling})
  \item[\texttt{[x]}] Audit query API contract specified
  \item[\texttt{[ ]}] Audit service implementation
        (\path{src/governance/audit_service/}) --- target: v1.0
  \item[\texttt{[ ]}] Immutable storage configured (append-only,
        hash chaining, external timestamp authority) --- target: v1.0
  \item[\texttt{[ ]}] Retention policy enforced (7 years) --- target: v1.0
\end{itemize}

\subsection{Pipeline integration rollout}
\label{sec:identity-integration-rollout}

The integration of the identity attribution audit service with the
three application pipelines (TB, Arabic AP, Simula) is sequenced
against explicit milestones, not an open-ended forward commitment.
Each pipeline has its own integration sub-ticket with an acceptance
criterion and a position in the dependency graph of
\cref{ch:reg-roadmap}; slippage on any row escalates to the GFAIC
governance committee. \textbf{Sequence, not dates}: the rows below
use the MoSCoW identifiers from \cref{tab:reg-milestones} so
implementation teams know what depends on what, not when a
calendar target is supposed to fire (that belongs in the steering
packet).

\begin{table}[htbp]
\centering
\caption{Audit service and pipeline integration rollout, keyed to
the MoSCoW priorities in \cref{tab:reg-milestones}. Pri.\ column:
M=Must (critical path), S=Should, C=Could.}
\label{tab:identity-rollout}
\footnotesize
\begin{tabularx}{\linewidth}{@{}l l l l X@{}}
\toprule
\textbf{ID} & \textbf{Prereq} & \textbf{Pri.} & \textbf{Ticket} & \textbf{Acceptance criterion} \\ \midrule
R1 & ---        & M & REG-AUDIT-001   & Audit service skeleton (reserve/complete endpoints, in-memory store): contract tests pass against the two endpoints; 404/409/500 paths exercised. \\
R2 & R1         & M & REG-AUDIT-002   & Immutable storage backend (append-only + hash chain + RFC~3161 timestamp): tamper-test suite passes; restored chain validates after adversarial edit. \\
R3 & R2         & M & REG-AUDIT-003   & 7-year retention + legal-hold API: retention policy enforced; legal-hold flag blocks GDPR erasure. \\
R4 & R3         & M & REG-INTG-TB-001 & TB dashboard integration: every dashboard decision \emph{and} AI commentary draft emits \texttt{reserve}/\texttt{complete}; 100\,\% coverage in \cref{ch:testing}. \\
R5 & R4         & S & REG-INTG-AP-001 & Arabic AP integration: OCR, field extraction, VAT rule evaluation each emit audit events; PII hashed. \\
R6 & R5         & S & REG-INTG-SIM-001 & Simula batch integration: training-data runs logged at batch granularity; per-example logging enabled under \texttt{AUDIT\_VERBOSE=1}. \\
R7 & R4         & M & GAIC-REV        & v1.0 compliance gate: quarterly GFAIC review confirms MAS MGF~2.1.2-002 compliance; traceability matrix \cref{tab:trace} flips from \texttt{partial} to \texttt{implemented} on this row. \\
\bottomrule
\end{tabularx}
\end{table}

\noindent
Until each pipeline's integration is verified, the governance status
for that pipeline's MAS MGF~2.1.2-002 compliance is \texttt{partial}
(see the traceability matrix in \cref{ch:refs}). Integration tickets
REG-INTG-TB-001 (R4) is a Must for the Regulations project and a Must
for TB; REG-INTG-AP-001 (R5) is a Should for Regulations and a Must
for Arabic AP (it is the \texttt{partial}-to-\texttt{compliant}
upgrade path for REG-MGF-2.1.2-002); REG-INTG-SIM-001 (R6) is Should
for Regulations and is not a v1.0 gate for Simula itself.
Table~\ref{tab:identity-rollout} is tracked in the programme status
dashboard and refreshed on every cycle.

\subsection{Compliance tests}

\begin{itemize}
  \item[\texttt{[ ]}] Per-pipeline compliance tests passing
        (\cref{ch:testing}, Section on audit coverage)
  \item[\texttt{[ ]}] Quarterly re-validation against the
        AIVerify subset (\cref{sec:aiverify-subset})
\end{itemize}